\Needspace{18\baselineskip}
\section{Sequential realization of simultaneous substitution}
\label{sec:sequential-expansion}
The simultaneous construction is now connected to the existing addressed
replacement operation. Every finite filling family can be performed by
ordinary shallow substitutions, and every completed schedule described below
has the same result up to explicit graph isomorphism. This supplies the
shallow compatibility obligation in Appendix~\ref{app:remaining-work}.
It does not identify these structural substitutions with UAM transitions,
nor by itself prove interchange for arbitrary deep sites; the latter is
established in Section~\ref{sec:deep-interchange}.

\subsection{Peeling one filling}
\begin{definition}[Residual filling family]\label{term:residual-family}
Fix a finite typed graph $C$ and a family $H=(H_f)_{f\in E_C}$ of
interface-compatible fillings. Select $e\in E_C$ and put $C'=C[e:=H_e]$,
using ordinary replacement. A retained edge of $C'$ is an original edge
$f\ne e$, and an inserted edge is an edge $x$ of $H_e$.
The residual family $R$ assigns $H_f$ to the former and the identity corolla
with label $\ell(x)$ to the latter. Thus subsequent expansion completes the
original assignments without replacing the newly inserted edges nontrivially.
\end{definition}
\lean{Graph/ExpansionPeel}{15}{19}{Graph.Net.ExpansionPeel.rest}

\begin{lemma}[Peeling isomorphism]
There is an explicit interface-preserving isomorphism
\[
             C'[R]\ \cong\ C[H].
\]
No acyclicity, absence of self-loops, or restriction on edge labels is needed.
\end{lemma}
\begin{proof}
First construct the bijection of internal vertices. An original host
vertex becomes its host copy. An internal vertex already inserted from
$H_e$ becomes the copy tagged by $e$ in $C[H]$. An internal vertex subsequently
inserted from $H_f$, with $f\ne e$, becomes the copy tagged by $f$.
A corolla has no internal vertices, so there is no fourth case.
The inverse retains host vertices and splits a filling tag $f$ according
to whether $f=e$. In the equality case, transporting along that equality
places the vertex in the already inserted $H_e$ carrier; otherwise it is
placed in the still-pending $H_f$ carrier. Case analysis proves both inverse
identities.
\lean{Graph/ExpansionPeel}{22}{26}{Graph.Net.ExpansionPeel.vertexTo}
\lean{Graph/ExpansionPeel}{28}{32}{Graph.Net.ExpansionPeel.vertexFrom}
\lean{Graph/ExpansionPeel}{34}{51}{Graph.Net.ExpansionPeel.vertices}

The edge bijection sends an edge of a pending $H_f$ to the pair $(f,x)$
and the sole edge of the corolla for an inserted $x\in E_{H_e}$ to $(e,x)$.
Its inverse again splits on equality with $e$. The only extra inverse
equation says that every element of $\mathrm{Fin}(1)$ is zero.
\lean{Graph/ExpansionPeel}{53}{55}{Graph.Net.ExpansionPeel.edgeTo}
\lean{Graph/ExpansionPeel}{57}{59}{Graph.Net.ExpansionPeel.edgeFrom}
\lean{Graph/ExpansionPeel}{61}{77}{Graph.Net.ExpansionPeel.edges}

It remains to verify incidence, rather than merely count vertices and edges.
The composite inclusion of an original host vertex agrees with its
simultaneous inclusion. The corresponding assertion for a vertex of $H_e$
holds directly for internal vertices; for boundary vertices it reduces to
the host assertion at the appropriate attachment port.
\lean{Graph/ExpansionPeel}{79}{81}{Graph.Net.ExpansionPeel.map_host}
\lean{Graph/ExpansionPeel}{83}{89}{Graph.Net.ExpansionPeel.map_filled}

For a retained edge $f$, each source or target port after the first
replacement is the host inclusion of its original port. This follows by
taking the indexed element of the mapped incidence list.
Consequently inclusion of any vertex of a pending $H_f$ commutes with the
vertex bijection: the internal case is definitional, and each boundary
case uses the appropriate retained-port identity followed by the host
identity.
\lean{Graph/ExpansionPeel}{91}{93}{Graph.Net.ExpansionPeel.sourcePort_retained}
\lean{Graph/ExpansionPeel}{95}{97}{Graph.Net.ExpansionPeel.targetPort_retained}
\lean{Graph/ExpansionPeel}{99}{113}{Graph.Net.ExpansionPeel.map_other}

For an inserted edge $x$ of $H_e$, the input of its residual corolla maps
to the simultaneous inclusion of the corresponding source port of $x$;
the output maps to the corresponding target port. These identities use
the already proved port formulas for an inserted edge of an ordinary plug,
followed by the filled-vertex identity above.
\lean{Graph/ExpansionPeel}{115}{120}{Graph.Net.ExpansionPeel.map_corolla_input}
\lean{Graph/ExpansionPeel}{122}{127}{Graph.Net.ExpansionPeel.map_corolla_output}

These bijections and incidence identities form the promised isomorphism.
Vertex and edge labels are preserved by inspection of the tags.
For a pending filling, preservation of its source and target lists follows
by composing list maps and applying the vertex identity pointwise.
For a corolla, its ordered list of inputs or outputs enumerates all port
indices. Mapping this list through the corresponding port function
recovers exactly the incidence list of $x$ in $H_e$.
This completes the proof of the lemma, including the ordering of incidence.
\end{proof}
\lean{Graph/ExpansionPeel}{131}{162}{Graph.Net.ExpansionPeel.iso}

\subsection{Stopping at arbitrary identity copies}
The termination condition must permit arbitrary graphs certified as
corolla copies, rather than require the literal corolla constructor.
The existing copy record specifies a sole edge with the original label
and the full ordered input and output lists.

\begin{lemma}[A corolla copy has no internal vertices]
If $H_e$ is a corolla copy, its internal carrier is empty.
\end{lemma}
\begin{proof}
If $v$ were internal, source coverage would place it in the source list of
some edge. That edge must be the sole edge of the copy, whose source list
contains only input vertices. The disjointness of input and internal
constructors gives a contradiction.
\end{proof}
\lean{Graph/ExpansionCopies}{13}{17}{Linear.ExpansionPaths.copy_no_internal}

\begin{lemma}[Restoring a host from copies]
If every $H_e$ is a corolla copy, then $C[H]\cong C$.
\end{lemma}
\begin{proof}
The internal-vertex bijection drops the host tag; its putative filling
cases are impossible by the preceding lemma. Its inverse inserts the host
tag. The edge bijection forgets the owning tag's unique filling edge,
sending $(e,x)$ to $e$; the inverse chooses the sole edge of the copy.
Uniqueness proves the inverse equations. Labels are the original labels
by the copy certificate. For each incidence list, expansion of the copy
gives the original host list mapped through host inclusion. The
internal-vertex bijection cancels that inclusion, including at boundary
vertices, so mapping back restores the original list.
\end{proof}
\lean{Graph/ExpansionCopies}{20}{27}{Linear.ExpansionPaths.copiesVertices}
\lean{Graph/ExpansionCopies}{30}{34}{Linear.ExpansionPaths.copiesEdges}
\lean{Graph/ExpansionCopies}{37}{62}{Linear.ExpansionPaths.copiesIso}

\subsection{Certified executions and order independence}
Fix a hypernet alphabet, a depth bound, and an outer interface.
\begin{definition}[Shallow substitution run]\label{term:shallow-run}
A shallow substitution run from $C$ to $D$ of length $n$ is either the
empty run ($n=0$, $D=C$), or an ordinary addressed replacement at a shallow
edge of $C$, followed by a run of length $n-1$ from the resulting graph.
Each step uses the existing operation $\operatorname{plugAt}$.
These are structural graph replacements, not machine evaluation steps.
\end{definition}
\lean{Graph/ExpansionSequential}{11}{18}{Linear.ShallowRun}

\begin{definition}[Certified filling execution]\label{term:filling-execution}
An execution is indexed by the current host $C$, its prescribed family
$H$, a length $n$, and a final graph $D$. It may stop with $n=0$ and $D=C$
when every assignment is certified as a corolla copy. Otherwise it selects
an edge $e$, performs $C[e:=H_e]$, and continues with the residual family
defined above. The length increases by one.
The definition records the actual steps and a legitimate stopping
condition; it does not assume any isomorphism of the final result.
\end{definition}
\lean{Graph/ExpansionSequential}{27}{32}{Linear.ExpansionSequential.Execution}

Every certified execution yields a shallow substitution run. Induction
sends the stopping constructor to the empty run, and each peeling
constructor to the corresponding addressed replacement followed by the
inductively obtained run.
\lean{Graph/ExpansionSequential}{35}{39}{Linear.ExpansionSequential.Execution.ordinary}

\begin{theorem}[Correctness of a completed execution]
An execution from $(C,H)$ to $D$ gives $D\cong C[H]$.
\end{theorem}
\begin{proof}
Induct on the execution. At termination, choose the certified corolla
copies and invert the host-restoration isomorphism.
At a step selecting $e$, the induction hypothesis identifies $D$ with
$C'[R]$. Compose this isomorphism with the peeling isomorphism
$C'[R]\cong C[H]$.
\end{proof}
\lean{Graph/ExpansionSequential}{43}{50}{Linear.ExpansionSequential.Execution.iso}

\begin{theorem}[Order independence]
Any two completed executions of the same prescribed family have
isomorphic final graphs, even if their lengths differ. They also
represent the same recursive graph class.
\end{theorem}
\begin{proof}
If their results are $D$ and $E$, correctness gives isomorphisms
$D\cong C[H]$ and $E\cong C[H]$. Compose the first with the inverse of
the second. An ordinary graph isomorphism lifts to a recursive graph
isomorphism, which implies equality in the recursive quotient.
This does not assert confluence for runs that change the prescribed
fillings, or for arbitrary UAM transitions.
\end{proof}
\lean{Graph/ExpansionSequential}{54}{59}{Linear.ExpansionSequential.Execution.order_independent}
\lean{Graph/ExpansionSequential}{62}{67}{Linear.ExpansionSequential.Execution.order_independent_modulo}

\subsection{Existence and a finite length bound}
\begin{definition}[Pending list]
A list $L$ of host edges is sufficient if every edge absent from it has
a corolla-copy assignment. Repetitions are allowed.
After selecting $e$, filter the tail of $L$: remove every occurrence of
$e$, and send each other edge $f$ to its retained copy in $C[e:=H_e]$.
Call the resulting list $L'$.
\end{definition}
\lean{Graph/ExpansionSequential}{71}{73}{Linear.ExpansionSequential.remaining}

Filtering cannot increase length. Furthermore, every retained edge whose
original lies in the tail appears in $L'$. The first claim is the standard
length bound for filtering with a partial map; for the second, the
original edge is a witness in the membership characterization of that map.
\lean{Graph/ExpansionSequential}{75}{76}{Linear.ExpansionSequential.remaining_length}
\lean{Graph/ExpansionSequential}{78}{81}{Linear.ExpansionSequential.remaining_mem}

Every edge outside $L'$ has a corolla-copy assignment in the residual
family. For a retained $f$, absence from $L'$ and the membership lemma
imply absence from the original pending list: $f\ne e$ excludes its head.
Use the original copy certificate. An inserted edge has a literal
corolla assignment, whose certificate uses its unique element of
$\mathrm{Fin}(1)$ and its defining incidence lists.
\lean{Graph/ExpansionSequential}{84}{109}{Linear.ExpansionSequential.remaining_copies}

\begin{theorem}[Finite sequential realization]
For every sufficient list $L$, there are $n\le |L|$ and a final graph $D$
with a certified execution of $(C,H)$ of length $n$ ending at $D$.
Consequently every finite family has an ordinary shallow substitution
run with final graph isomorphic to $C[H]$. Its length is at most the
length of the host's finite edge enumeration.
\end{theorem}
\begin{proof}
Use strong induction on the pending-list length, allowing both the host
and family to vary. An empty list certifies every assignment as an
identity copy, so stop. For $e::L$, first replace $e$ with $H_e$.
The residual list has length at most $|L|$, strictly less than the original
length, and the preceding lemma supplies the required certificates
outside it. Apply the induction hypothesis and prepend the selected
step. Its total length is at most $|L|+1$.
For an arbitrary finite family, use the host's complete edge enumeration;
there are no edges outside it. Forget the resulting execution to obtain
an ordinary run, and apply correctness. The enumeration may contain
repetitions, so the stated bound is its list length, not an asserted
minimal number of substitutions.
\end{proof}
\lean{Graph/ExpansionSequential}{114}{142}{Linear.ExpansionSequential.realize_list}
\lean{Graph/ExpansionSequential}{146}{150}{Linear.ExpansionSequential.realize}

The same witness supplies equality of the final result and simultaneous
expansion in the recursive graph quotient: lift the constructed graph
isomorphism and use quotient soundness.
\lean{Graph/ExpansionSequential}{153}{156}{Linear.ExpansionSequential.realize_modulo}

\subsection{Two-hole examples and scope}
Both filling families in the alternating one-way example admit ordinary
runs of length at most two. Apply finite realization to each family;
the explicit context enumeration has length two.
\lean{Graph/SequentialExamples}{10}{17}{Linear.AlternatingFillings.sequential_realizations}

For any filling family on this context, there are completed executions
starting by replacing hole $0$ and by replacing hole $1$. Each requires
at most one further step, and their endpoints are isomorphic. To prove
this, verify that the two-element list contains every context edge.
After either initial choice, the residual-list certificate permits
finite realization on the remaining list of length at most one.
Prepend the chosen step to each execution and apply order independence.
The statement retains the residual execution witnesses after each
distinct first choice, so it records which hole was filled first.
\lean{Graph/SequentialExamples}{21}{54}{Linear.AlternatingFillings.reverse_schedule}

The dependency chain is now: explicit peeling and identity-copy
isomorphisms, then correctness of executions, then finite realization
and schedule independence. These results remove a representation gap
between simultaneous shallow specimens and ordinary graph replacement.
Section~\ref{sec:deep-interchange} adds a quotient of host/site pairs and
independent deep-site interchange. Deep specimen substitution and transport
of the entire focused-state and UAM structure along arbitrary deep replacement
schedules remain separate obligations.
